%%%%%%%%%%%%%%%%%%%%%%%%%%%%%%%%%%%%%%%%%%%%%%%%%%%%%
%
%           EE474 Hardware and Basic I/O
%
   % combined into notes pack	%<stream, c>

\documentclass{article}	%<nh>

\usepackage{graphicx}	%<shn>
\usepackage{hyperref}	%<shn>
\usepackage{listings}	%<shn>

	%<*hn>
\oddsidemargin 0in
\evensidemargin 0in
\textwidth 6.5in
\textheight 8.7in
\topmargin 0.2in


	%<*shn>
\linespread{1.0}


\begin{document}	%<shn>

	%<*>



\section{Basic Machine Organization}
\subsection{Registers}

Registers are an array of D flip-flops which can store a collection of bits.

An n bit register has n inputs, n outputs, and one clock line.

\begin{center}
   \includegraphics[height=2.5in]{hwio_figs/00156.eps}   \\	%<chn>

\end{center}


\subsection{Busses}

A bus is a parallel bi-directional data path

Multiple devices can send and receive on one bus

\begin{center}
   \includegraphics[height=2.5in]{hwio_figs/00155.eps}   \\	%<chn>
\end{center}

Address: N bits which specify {\it which location}.

Data:   Contents to/from memory or I/O device

Control:  'traffic signals'



More detailed view:

Address Comparator

\begin{itemize}
 \item A combinatorial logic circuit at each device
 \item 2 $n$-bit binary inputs:
   \begin{enumerate}
          \item ADDR bus (changing)
          \item Device address (fixed)
    \end{enumerate}

 \item Control input
 \item Output = 1 iff input1 == input2
\end{itemize}


Control bus is gated by ADDR comparator

Device register is strobed by CTL pulse only if addresses match.

\begin{center}
   \includegraphics[height=2.5in]{hwio_figs/00157.eps}   \\	%<chn>
\end{center}




Bus Logic Notation:

Below is a simple notation for showing timing relationships of signals on
busses.

\begin{center}
   \includegraphics[height=2.5in]{hwio_figs/00158.eps}   \\	%<chn>

\end{center}

Problem:  How can multiple logic outputs be connected to the
same wire?

Solution:  Tri-state logic.

\begin{center}
   \includegraphics[height=2.0in]{hwio_figs/00159.eps}   \\	%<chn>
\end{center}
States:  { Logic 1, Logic 0, switch OFF }

CPU and all devices can ``talk'' on data bus.

Only {\bf ONE} device my have the tri-state switch closed at
any time!


\subsection{Bus Timing}

Write Cycle

\begin{center}
   \includegraphics[height=2.5in]{hwio_figs/00160.eps}   \\	%<chn>

\end{center}


Read Cycle

\begin{center}
   \includegraphics[height=2.5in]{hwio_figs/00161.eps}   \\	%<chn>
\end{center}





\subsection{Microprocessor: Register View}

\begin{center}
   \includegraphics[height=2.5in]{hwio_figs/00162.eps}   \\	%<chn>

\end{center}

% 

\begin{itemize}
\item ALU: Arithmetic \& Logic Unit
\item PC:  Program Counter. Holds Address of next instruction.
\item IR:  Instruction Register. Holds current instruction.
\item ADDR: Address Register. Holds address of next bus access.
\item GP\#: General Purpose Registers.  Hold intermediate results.
\item A1, A2: ALU Arguments. Hold inputs to an ALU operation.
\item r: 'Result', holds result of ALU operation.
\item FLAGS: a set of bits which tell things like zero/non-zero, negative, etc about the last Result.
\end{itemize}



\subsection{Steps to add two numbers}

 C statement: {\tt c = a + b ;}
\begin{enumerate}

\item  Load addr of {\tt a} into ADDR
\item  Wait for data from memory
\item  Clock data into A1
\item  Load addr of {\tt b} into ADDR
\item  Wait for data from memory
\item  Clock data into A2
\item  Send ADD command to ALU
\item  Load addr of {\tt c}
\item  Transfer RES to data bus
\item  Wait for data write to memory.

\end{enumerate}


\section{Machine Language Instructions}

The above example may be controlled by one or more machine instructions.
A machine instruction is binary data typically broken up into fields:

The Operation Code (Op-Code)

One, two, or three Operands.

Each instruction is typically between one and eight bytes.

Examples:
\begin{center}
   \includegraphics[height=2.5in]{hwio_figs/00163.eps}   \\	%<chn>

\end{center}


\subsection{Example: Add two numbers}
Note:  All assembler below is pseudo-code!

C: {\tt c = a + b;}

Assembler:

\begin{tabular}{|r|l|l|}
\hline
mem addr        & instruction   & comment\\
\hline
0xA000          & MOV  a, A1   & move mem location a to ALU Arg 1 \\
\hline
0xA002          & MOV  b, A2    & move mem location b to ALU Arg 2 \\
\hline
0xA004          & ADD           & a one-byte instruction \\
\hline
0xA005          & MOV  r, c     & move ALU result to mem location c\\
\hline
\end{tabular}


There are many more types of operands in most architectures.  Examples:

\begin{itemize}
\item INDEXED
\begin{verbatim}
MOV   0xCA04(R2), G1    \\move data from memory location
                        \\ 0xCA04 + contents of R2 to register G1
\end{verbatim}

\item REGISTER INDIRECT
\begin{verbatim}
MOV   (R2), G1          \\move data from memory location
                        \\ contents of R2
\end{verbatim}

\end{itemize}




\section{Basic I/O programming}

{\bf Intel} I/O instructions

{\tt IN         R, char }
where {\tt char} is an 8-bit constant

{\tt OUT        char, R }

In intel architecture, {\tt R} can be {\tt AL, AX, or EAX} registers.




{\bf Atmel} I/O Instructions 

Example:   AtMega2560 Chip (Arduino Mega)

\begin{itemize}

    \item ``Memory Mapped" I/O
    \item Devices and Memory {\bf share } the same address space
    \item Device registers are just like memory locations / variables. 
    \item Arduino libraries predeclare correct memory address for each register you need. 
\end{itemize}
\end{document}	%<shn>

